% Zone „Das kennst du schon“ – aus bank/potenzen-wurzeln/zone.jsonl (zusammenbau v0.1)
\einheitenkopf[zone][Kennst du schon]{Das kennst du schon}
\setcounter{aufgabe}{0}
%% TODO Grundvorstellungs-Aufgabe als erste Hauptnummer der Zone (2.8) fehlt in der Bank

%% TODO Titel: Ich-kann-Satz fehlt in der Bank (Platzhalter: Kettenname) – Nr. 1
%% TODO Anweisung: ein Satz über der ersten Teilaufgabe fehlt in der Bank – Nr. 1
\begin{aufgabe}{Kleines Einmaleins und Malnehmen mehrstelliger Zahlen im Kopf oder halbschriftlich, Malketten mit gleichen Faktoren}
%% TODO Darstellung für schwach fehlt in der Bank (potenzen-wurzeln-zone-f1-v1; Abschnitt „Für schwache Schüler“)
\swz{$7 \cdot 8$ – wie viel?}{}
%% TODO Darstellung für schwach fehlt in der Bank (potenzen-wurzeln-zone-f1-v3; Abschnitt „Für schwache Schüler“)
\swz{$23 \cdot 4$ – wie viel?}{}
\end{aufgabe}

%% TODO Titel: Ich-kann-Satz fehlt in der Bank (Platzhalter: Kettenname) – Nr. 2
%% TODO Anweisung: ein Satz über der ersten Teilaufgabe fehlt in der Bank – Nr. 2
\begin{aufgabe}{Quadratzahlen bis 100 aus dem Kopf, Kubikzahlen zwei hoch drei, drei hoch drei, zehn hoch drei}
%% TODO Darstellung für schwach fehlt in der Bank (potenzen-wurzeln-zone-f2-v1; Abschnitt „Für schwache Schüler“)
\swz{$7 \cdot 7$ – wie viel?}{}
%% TODO Darstellung für schwach fehlt in der Bank (potenzen-wurzeln-zone-f2-v3; Abschnitt „Für schwache Schüler“)
\swz{$3 \cdot 3 \cdot 3$ – wie viel?}{}
\end{aufgabe}

%% TODO Titel: Ich-kann-Satz fehlt in der Bank (Platzhalter: Kettenname) – Nr. 3
%% TODO Anweisung: ein Satz über der ersten Teilaufgabe fehlt in der Bank – Nr. 3
\begin{aufgabe}{Vorzeichenregel beim Malnehmen (minus mal minus gibt plus), Anzahl der Minuszeichen zählen, Quadrat einer negativen Zahl mit und ohne Klammer}
%% TODO Darstellung für schwach fehlt in der Bank (potenzen-wurzeln-zone-f3-v1; Abschnitt „Für schwache Schüler“)
\swz{$(-4) \cdot (-5)$ – wie viel?}{}
%% TODO Darstellung für schwach fehlt in der Bank (potenzen-wurzeln-zone-f3-v3; Abschnitt „Für schwache Schüler“)
\swz{$(-2) \cdot (-3) \cdot (-5)$ – wie viel?}{}
\end{aufgabe}

%% TODO Titel: Ich-kann-Satz fehlt in der Bank (Platzhalter: Kettenname) – Nr. 4
%% TODO Anweisung: ein Satz über der ersten Teilaufgabe fehlt in der Bank – Nr. 4
\begin{aufgabe}{Dezimalzahlen multiplizieren (Kommastellen zählen), Brüche multiplizieren}
%% TODO Darstellung für schwach fehlt in der Bank (potenzen-wurzeln-zone-f4-v1; Abschnitt „Für schwache Schüler“)
\swz{$0{,}5 \cdot 6$ – wie viel?}{}
%% TODO Darstellung für schwach fehlt in der Bank (potenzen-wurzeln-zone-f4-v3; Abschnitt „Für schwache Schüler“)
\swz{$1{,}2 \cdot 0{,}3$ – wie viel?}{}
\end{aufgabe}

%% TODO Titel: Ich-kann-Satz fehlt in der Bank (Platzhalter: Kettenname) – Nr. 5
%% TODO Anweisung: ein Satz über der ersten Teilaufgabe fehlt in der Bank – Nr. 5
\begin{aufgabe}{Dezimalzahlen und Brüche vergleichen und ordnen, Bruch in Dezimalzahl umwandeln, Prozent in Dezimalzahl, negative Zahlen auf der Zahlengerade ordnen}
%% TODO Darstellung für schwach fehlt in der Bank (potenzen-wurzeln-zone-f5-v1; Abschnitt „Für schwache Schüler“)
\swz{$0{,}9$ oder $0{,}4$ – welche Zahl ist größer?}{}
%% TODO Darstellung für schwach fehlt in der Bank (potenzen-wurzeln-zone-f5-v3; Abschnitt „Für schwache Schüler“)
\swz{$35\,\%$ als Dezimalzahl – wie viel?}{}
\end{aufgabe}

%% TODO Titel: Ich-kann-Satz fehlt in der Bank (Platzhalter: Kettenname) – Nr. 6
%% TODO Anweisung: ein Satz über der ersten Teilaufgabe fehlt in der Bank – Nr. 6
\begin{aufgabe}{Taschenrechner}
%% TODO Darstellung für schwach fehlt in der Bank (potenzen-wurzeln-zone-f6-v1; Abschnitt „Für schwache Schüler“)
\swz{$43 \cdot 38$ – wie viel? Rechne mit dem Taschenrechner.}{}
%% TODO Darstellung für schwach fehlt in der Bank (potenzen-wurzeln-zone-f6-v3; Abschnitt „Für schwache Schüler“)
\swz{$(3{,}8 + 5{,}7) \cdot 2{,}4$ – wie viel? Rechne mit dem Taschenrechner.}{}
\end{aufgabe}

%% TODO Titel: Ich-kann-Satz fehlt in der Bank (Platzhalter: Kettenname) – Nr. 7
%% TODO Anweisung: ein Satz über der ersten Teilaufgabe fehlt in der Bank – Nr. 7
\begin{aufgabe}{Große Zahlen lesen und schreiben (Dreiergruppen, Zahlwörter bis Billion), Stellenwerttafel, mal zehn, hundert und tausend durch Nullen anhängen}
%% TODO Darstellung für schwach fehlt in der Bank (potenzen-wurzeln-zone-f7-v1; Abschnitt „Für schwache Schüler“)
\swz{$38 \cdot 100$ – wie viel?}{}
%% TODO Darstellung für schwach fehlt in der Bank (potenzen-wurzeln-zone-f7-v3; Abschnitt „Für schwache Schüler“)
\swz{Drei Millionen vierhunderttausend – in Ziffern?}{}
\end{aufgabe}

%% TODO Titel: Ich-kann-Satz fehlt in der Bank (Platzhalter: Kettenname) – Nr. 8
%% TODO Anweisung: ein Satz über der ersten Teilaufgabe fehlt in der Bank – Nr. 8
\begin{aufgabe}{Kommaverschiebung}
%% TODO Darstellung für schwach fehlt in der Bank (potenzen-wurzeln-zone-f8-v1; Abschnitt „Für schwache Schüler“)
\swz{$3{,}75 \cdot 10$ – wie viel?}{}
%% TODO Darstellung für schwach fehlt in der Bank (potenzen-wurzeln-zone-f8-v3; Abschnitt „Für schwache Schüler“)
\swz{$0{,}062 \cdot 1\,000$ – wie viel?}{}
\end{aufgabe}

%% TODO Titel: Ich-kann-Satz fehlt in der Bank (Platzhalter: Kettenname) – Nr. 9
%% TODO Anweisung: ein Satz über der ersten Teilaufgabe fehlt in der Bank – Nr. 9
\begin{aufgabe}{Runden von Dezimalzahlen auf eine und zwei Stellen, Näherungswert mit ≈ schreiben}
%% TODO Darstellung für schwach fehlt in der Bank (potenzen-wurzeln-zone-f9-v1; Abschnitt „Für schwache Schüler“)
\swz{$4{,}372$ – auf eine Stelle nach dem Komma gerundet?}{}
%% TODO Darstellung für schwach fehlt in der Bank (potenzen-wurzeln-zone-f9-v3; Abschnitt „Für schwache Schüler“)
\swz{$0{,}0648$ – auf zwei Stellen nach dem Komma gerundet?}{}
\end{aufgabe}

%% TODO Titel: Ich-kann-Satz fehlt in der Bank (Platzhalter: Kettenname) – Nr. 10
%% TODO Anweisung: ein Satz über der ersten Teilaufgabe fehlt in der Bank – Nr. 10
\begin{aufgabe}{Vorzeichenregel beim Malnehmen (minus mal minus gibt plus), Anzahl der Minuszeichen zählen, Quadrat einer negativen Zahl mit und ohne Klammer – Fehler finden}
\swz[3]{Wo ist der Fehler? $(-11)^2 = -121$}{}
\end{aufgabe}

%% TODO Titel: Ich-kann-Satz fehlt in der Bank (Platzhalter: Kettenname) – Nr. 11
%% TODO Anweisung: ein Satz über der ersten Teilaufgabe fehlt in der Bank – Nr. 11
\begin{aufgabe}{Vorzeichenregel beim Malnehmen (minus mal minus gibt plus), Anzahl der Minuszeichen zählen, Quadrat einer negativen Zahl mit und ohne Klammer – selbst rechnen}
%% TODO Darstellung für schwach fehlt in der Bank (potenzen-wurzeln-zone-f3-v7; Abschnitt „Für schwache Schüler“)
\swz{$(-15)^2$ – wie viel?}{}
\end{aufgabe}
